\begin{frame}{Initialize a project with an AI agent}
\framesubtitle{Manual Section 1.9: Initializing a New Project from This Template}
{\small Open the project copy in your AI tool, attach the available project evidence, and replace \mbox{\texttt{[Project Name]}}. Enter:\par}
\vspace{0.2cm}
\begingroup\small\ttfamily\raggedright
\noindent Read dsnAIAgentGuide.tex and use the AIS-02 task steps for [Project Name]. Execute using Appendix D. Review the supplied project evidence and create a gap list before editing. Identify missing scope, stakeholders, interfaces, constraints, quality attributes, risks, evidence, and approvals. Identify the available domain profiles and ask me to confirm whether each is Applicable, Not Applicable, or TBD. Read each Applicable or relevant TBD profile and add its concerns to the gap list. After the gap list is reviewed, replace only supported placeholders and develop the document in dependency order. Compile the complete PDF and report unresolved human decisions.\par
\endgroup
\vspace{0.2cm}
{\small\textbf{Be ready to supply:} the project namespace, metadata, profile decisions, and missing evidence.}
\end{frame}
\note{Sources: Introduction Section 1.9 (Section:InitializeNewProject) and AI Agent Guide, Initialize the Architecture Book (Section:AIInitializePrompt). The displayed request reproduces the manual's initialization prompt. Start with a project copy of the template and provide the AI tool access to that workspace and supplied evidence. Appendix D is the AI Implementation Workflow. AIS-02 requires a human-confirmed namespace before creating or renaming files and recorded profile applicability decisions. Review the gap list before detailed work; unsupported facts remain gaps.}
